% !TEX root = ../main.tex
\documentclass[../main.tex]{subfiles}

\begin{document}

\chapter{Synthesis and Discussion}
\labch{synthesis}
\margintoc

\begin{chapteroverview}
This chapter synthesizes findings across all empirical chapters and discusses their broader implications. We integrate insights from multimodal cartography to articulate how role constellations emerge, stabilize, and transform in sustainability transitions. The chapter develops a unified theoretical framework, addresses methodological contributions and limitations, and identifies practical implications for transition governance.
\end{chapteroverview}

%===============================================================================
\section{Introduction}
\labsec{synthesis-intro}
%===============================================================================

The fundamental question driving this dissertation is: \emph{What are the role constellation patterns for each EIT actor, as revealed through each modality?}\sidenote{\iconidea\ \textbf{Integration challenge}---Each modality reveals different facets of the same organizations. Synthesis requires both combining and contrasting these views.} The preceding empirical chapters answered this question from different perspectives:

\begin{itemize}[leftmargin=*, itemsep=0.3em]
\item Chapter~\ref{ch:interstitial} revealed \textbf{claimed roles}---what organizations say about themselves through website positioning and vocabulary
\item Chapter~\ref{ch:rolefield} revealed \textbf{relational roles}---how organizations position themselves in social media networks and discourse
\item Chapter~\ref{ch:venture} revealed \textbf{resourced roles}---what portfolio compositions and investment patterns reveal about enacted functions
\item Chapter~\ref{ch:visual} revealed \textbf{embodied roles}---how visual registers construct and communicate institutional identity
\item Chapter~\ref{ch:power} revealed \textbf{power dynamics}---how discourse constructs agency hierarchies among ecosystem actors
\end{itemize}

Each modality uncovered distinct role constellation patterns across the EIT ecosystem. But the question remains: \emph{How do these views combine to reveal the complete role profile of each intermediary actor?}

This synthesis chapter addresses three integrative questions:

\begin{enumerate}
    \item \textbf{Integration:} How do signals from different modalities converge or diverge, and what do these patterns tell us about role dynamics?
    \item \textbf{Theory:} What theoretical contributions emerge from multimodal role constellation cartography?
    \item \textbf{Practice:} What are the implications for transition governance and innovation policy?
\end{enumerate}

%===============================================================================
\section{Cross-Modal Triangulation}[Triangulation]
\labsec{triangulation}
%===============================================================================

The central methodological challenge of multimodal cartography is integration: how do we synthesize signals from five different data modalities into coherent role profiles?\sidenote{\textbf{Triangulation logic}---Classic triangulation seeks convergence \cite{denzin1970triangulation,jick1979mixing}. Multimodal cartography goes further: divergence is equally informative.} Our approach treats both convergence and divergence as informative. When multiple modalities agree on an actor's role profile, we gain confidence in the attribution. When modalities disagree, we have not encountered noise---we have found something analytically interesting: aspiration gaps (claimed roles exceeding enacted roles), perception gaps (attributed roles differing from claimed roles), resource gaps (claimed roles lacking investment support), or strategic gaps (deliberate divergence for positioning purposes).

\subsection{Role Spheres: The X-Curve as Conceptual Lens}

The X-curve framework \parencite{hebinck2022xcurve} identifies eight transition phases organized along two arms: \textit{build-up} (experimentation, acceleration, emergence, institutionalization) and \textit{breakdown} (destabilization, chaos, phase-out, decline). We propose treating each phase as a \textbf{Role Sphere}---a spatio-temporal site where particular role performances cluster.\sidenote{\textbf{Role Sphere}---A region of the transition landscape characterized by a distinct mix of role demands, temporal dynamics, and relational structures.} Each sphere has distinctive spatial position (where on the X-curve), temporal character (pace and direction of change), and relational structure (how actors connect).

Crucially, the Multi-Level Perspective \parencite{geels2002technological,geels2004process} suggests these spheres are \textbf{coupled}, not independent. Build-up activities can trigger breakdown: niche innovations that scale successfully destabilize incumbent regimes. But the coupling runs in both directions. Build-up can also \textbf{crowd out} breakdown---absorbing the resources, attention, and institutional energy that deliberate phase-out would require.\sidenote{\textbf{The crowding-out mechanism}---If all ecosystem funding, talent, and political capital flows toward innovation, nothing remains for managing decline.} An intermediary ecosystem saturated with experimentation and acceleration leaves no institutional space for actors focused on managed decline. Moreover, build-up generates its own lock-ins: legacy investments create sunk costs that resist reallocation \parencite{unruh2000understanding}; institutional commitments produce path dependencies; and the political economy of innovation funding rewards scaling over sunsetting. This echoes broader findings on the systematic neglect of subtractive solutions \parencite{adams2021subtractive}.

\paragraph{What the pipelines found.} Applying this lens across all five modalities produces a stark asymmetry. Every pipeline detected abundant build-up sphere signatures and \textbf{zero breakdown sphere signatures}:\sidenote{\textbf{Cross-modal convergence on absence}---The five modalities diverge productively on \textit{which} build-up roles organizations occupy. But they converge completely on \textit{which Role Spheres} those organizations inhabit: build-up only.}

\begin{itemize}[leftmargin=*, itemsep=0.3em]
\item \textbf{Website vocabularies} cluster into five sociotypes (\cref{ch:interstitial}), all oriented toward creation, connection, and growth. No organization's self-presentation includes phase-out or managed decline language.
\item \textbf{Social media networks} organize into discourse sub-regions (\cref{ch:rolefield}), all structured around build-up activities. No cluster centers on exit or wind-down discourse.
\item \textbf{Investment patterns} reveal portfolio configurations spanning acceleration, scaling, and market entry (\cref{ch:venture}). No funding instrument targets deliberate phase-out or exnovation support.
\item \textbf{Visual registers} deploy three gazes---institutional, innovation, collaborative (\cref{ch:visual})---all projecting forward-looking dynamism. No visual register codes for sunset or transition-away imagery.
\item \textbf{Discourse networks} construct high-agency verb patterns around launching, developing, and supporting (\cref{ch:power}). No actor-action-target triple encodes phase-out agency.
\end{itemize}

The cross-modal convergence is total---not on \textit{what} build-up roles organizations occupy (there the modalities diverge productively), but on \textit{which Role Spheres they inhabit}. Every modality independently confirms exclusive occupation of build-up spheres. This raises the crowding-out question directly: does the density of build-up activity itself suppress the emergence of breakdown roles? If intermediaries channel all institutional energy into experimentation, acceleration, and scaling, the phase-out functions that exnovation scholarship identifies as essential \parencite{david2017exnovation,davidson2019exnovating} may never find organizational carriers.

%===============================================================================
\section{Cross-Chapter Integration}[Integration]
\labsec{cross-chapter}
%===============================================================================

Comparing findings across the empirical chapters reveals systematic patterns in how role constellations manifest across intermediation domains.\sidenote{\icondata\ \textbf{8 KICs, 5 modalities}---Cross-chapter integration draws on 40 KIC-modality combinations, enabling robust pattern identification.}

\subsection{Cross-Domain Patterns}

Despite domain differences, several patterns recur:

\begin{enumerate}
\item \textbf{The knowledge triangle imbalance persists.} Across all domains, business vocabulary dominates over research and education orientations (see \cref{ch:interstitial}).\sidenote{\textbf{EIT's founding vision}---The 2008 EIT Regulation envisioned equal integration of business, research, and education. Reality diverges.} The EIT's founding vision of balanced integration remains unrealized in organizational self-presentation.

\item \textbf{Connector inflation is universal.} Every domain shows more claimed connector roles than enacted bridging positions \cite{howells2006intermediation,kivimaa2019innovation}. Intermediation is aspirationally attractive but structurally scarce.

\item \textbf{Sociotypes cluster similarly.} The five sociotype configurations (Entrepreneurial, Research, Connective, Public, Accountability) emerge across domains, suggesting these are general features of innovation intermediary ecosystems rather than domain-specific artifacts.

\item \textbf{Peripheral actors remain invisible.} Multimodal cartography, like any method, privileges visible actors. Small organizations, failed ventures, and actors outside digital traces remain unmapped.
\end{enumerate}


\subsection{The Breakdown Gap: Cross-Modal Absence}

Perhaps the most striking cross-modal finding is what does \emph{not} appear: \textbf{breakdown and phase-out roles are completely absent across all five modalities}.\sidenote{\textbf{Systematic absence}---Not a gap in one data source but consistent absence across claimed, relational, resourced, embodied, and discursive dimensions.} The X-curve framework \parencite{hebinck2022xcurve} demonstrates that sustainability transitions require both build-up and breakdown dynamics---what Turnheim and Geels call ``regime destabilisation'' \parencite{turnheim2012destabilisation,turnheim2015regime}. Exnovation scholarship emphasizes that phase-out requires deliberate governance \parencite{david2017exnovation,davidson2019exnovating,heyen2017phaseout}. Yet our multimodal cartography reveals:

\begin{center}
\small
\begin{tabular}{@{}p{2.8cm}p{3.5cm}p{4.5cm}@{}}
\toprule
\textbf{Modality} & \textbf{Build-Up Presence} & \textbf{Breakdown Presence} \\
\midrule
Websites (claimed) & 5 sociotypes, all build-up & Zero phase-out claims \\
Social Media (relational) & 10 sub-regions, all build-up & Zero breakdown clusters \\
Crunchbase (resourced) & 398K organizations, all growth & Zero managed-decline focus \\
Visual (embodied) & 3 gazes, all creation & Zero sunset/exit imagery \\
News (discourse) & High-power build-up verbs & Zero phase-out agency \\
\bottomrule
\end{tabular}
\captionof{table}{Cross-modal absence of breakdown/phase-out roles}
\label{tab:phaseout-gap}
\end{center}

This finding has three possible interpretations. \textbf{Strategic silence}: organizations may perform phase-out support but avoid claiming it due to political sensitivity---yet the absence extends to \emph{enacted} behavior, not just self-presentation. \textbf{Structural gap}: the intermediary ecosystem genuinely lacks actors focused on breakdown dynamics; exnovation intermediaries may need to be purpose-built with distinct mandates \parencite{rosenbloom2020deliberate}. \textbf{Categorical impossibility}: the concept of ``innovation intermediary'' may be structurally incompatible with breakdown roles---one cannot be an \emph{innovation} intermediary while facilitating \emph{exnovation}, echoing Schumpeter's insight that creation and destruction involve fundamentally different logics \parencite{schumpeter1942capitalism}.

Regardless of interpretation, the policy implication is clear: \textbf{if sustainability transitions require managed decline alongside scaling alternatives, the existing intermediary ecosystem is structurally incomplete}. Policymakers cannot assume that innovation intermediaries will spontaneously extend into phase-out functions \parencite{johnstone2020destabilisation}. Dedicated exnovation intermediaries, explicit phase-out mandates, or distinct funding instruments may be necessary \parencite{rinscheid2021phaseout}.

This finding connects to the theoretical framework developed in \cref{ch:cas}. Complex adaptive systems require functional diversity for resilience. An ecosystem with actors oriented only toward build-up lacks the capacity for managed decline---making transitions more disruptive than necessary when regime change eventually occurs.

%===============================================================================
\section{Theoretical Contributions}[Theory]
\labsec{theoretical-contributions}
%===============================================================================

This dissertation makes several theoretical contributions to the study of roles in sustainability transitions. We begin by evaluating the fifteen CAS propositions developed in \cref{ch:cas} against the empirical findings of \cref{ch:interstitial,ch:rolefield,ch:venture,ch:visual,ch:power}.

\subsection{Proposition Evaluation}

\cref{ch:cas} derived fifteen propositions from five CAS mechanisms. \cref{tab:proposition-eval} evaluates each against the empirical evidence.\sidenote{\textbf{Evaluation protocol}---``Supported'' indicates consistent cross-modal evidence; ``Partially supported'' indicates evidence from some modalities or KICs; ``Not directly assessed'' indicates the proposition concerns dynamics not captured by our cross-sectional design.}

\begin{center}
\small
\tableminimal
\begin{tabular}{@{}p{0.4cm}p{4.8cm}p{1.5cm}p{5cm}@{}}
\toprule
& \minimalhead{Proposition} & \minimalhead{Status} & \minimalhead{Evidence} \\
\midrule
\multicolumn{4}{l}{\textit{Emergence}} \\
\addlinespace
P1 & Role multiplicity is the norm & Supported & All five sociotypes contain organizations performing multiple roles; single-role actors are rare across KICs (\cref{ch:interstitial}) \\
\addlinespace
P2 & Constellation functions exceed sum of parts & Supported & Cross-modal triangulation reveals system-level patterns (connector gap, breakdown gap) invisible at actor level (\cref{ch:interstitial,ch:rolefield}) \\
\addlinespace
P3 & Emergence accelerates with interaction density & Partially & Older KICs (Climate, Digital) show more differentiated role structures, consistent with longer interaction histories (\cref{ch:venture}) \\
\midrule
\multicolumn{4}{l}{\textit{Self-Organization}} \\
\addlinespace
P4 & Role differentiation without central coordination & Supported & Five sociotypes emerge across KICs despite different mandates; structure reflects self-organization, not policy design (\cref{ch:interstitial}) \\
\addlinespace
P5 & Hierarchy from accumulated advantage & Supported & Network centrality correlates with resource access and longevity, not formal status (\cref{ch:rolefield,ch:venture}) \\
\addlinespace
P6 & Niche formation reduces competition & Partially & Sociotype differentiation is evident, but connector inflation suggests incomplete niche separation (\cref{ch:interstitial,ch:rolefield}) \\
\midrule
\multicolumn{4}{l}{\textit{Co-Evolution}} \\
\addlinespace
P7 & Constellation--policy co-evolution & Not assessed & Cross-sectional design cannot capture temporal co-evolution \\
\addlinespace
P8 & Success alters fitness landscape for others & Partially & KIC-level analysis shows that dominant actors reshape available positioning for others (\cref{ch:venture,ch:power}) \\
\addlinespace
P9 & Co-evolution with technology trajectories & Not assessed & Technology maturation dynamics require longitudinal data \\
\midrule
\multicolumn{4}{l}{\textit{Non-Linearity}} \\
\addlinespace
P10 & Punctuated equilibrium patterns & Not assessed & Temporal dynamics beyond cross-sectional snapshot \\
\addlinespace
P11 & Small interventions, disproportionate effects & Not assessed & Would require intervention studies or longitudinal natural experiments \\
\addlinespace
P12 & Path dependence constrains repertoires & Partially & Founding-era KICs show more entrenched role profiles than newer ones (\cref{ch:venture}) \\
\midrule
\multicolumn{4}{l}{\textit{Feedback Dynamics}} \\
\addlinespace
P13 & Reinforcing feedback produces concentration & Supported & Skewed degree distributions in all network modalities (\cref{ch:rolefield,ch:power}) \\
\addlinespace
P14 & Balancing feedback maintains diversity & Partially & Five sociotypes persist across KICs, but breakdown gap suggests balancing mechanisms are insufficient \\
\addlinespace
P15 & Fitness requires multi-dimensional adaptation & Supported & Organizations visible in multiple modalities occupy more central positions (\cref{ch:rolefield,ch:venture}) \\
\bottomrule
\end{tabular}
\captionof{table}{Evaluation of fifteen CAS propositions (\cref{ch:cas}) against empirical evidence from \cref{ch:interstitial,ch:rolefield,ch:venture,ch:visual,ch:power}. ``Not assessed'' propositions require longitudinal designs beyond this dissertation's cross-sectional scope.}
\labtab{proposition-eval}
\end{center}

Of the fifteen propositions, seven receive support from cross-modal evidence, four are partially supported, and four concern temporal dynamics that our cross-sectional design cannot assess directly. The unassessed propositions (P7, P9, P10, P11) share a common requirement: longitudinal data capturing constellation evolution over time. This represents the most significant empirical frontier for future research (see \cref{ch:conclusion}).

\subsection{Role Constellations as Complex Adaptive Systems}

The CAS reconceptualization (\cref{ch:cas}) provides mechanisms---emergence, self-organization, co-evolution, non-linearity, feedback loops---that explain observed constellation dynamics while enabling computational operationalization \sidecite{holland1995hidden,kauffman1993origins}.\sidenote{\textbf{Theory-method bridge}---CAS mechanisms are simultaneously theoretically meaningful and computationally tractable---enabling the bridge this dissertation builds.} Key theoretical insights include:

\begin{itemize}[leftmargin=*, itemsep=0.3em]
\item \textbf{Emergence is real.} Constellation-level properties cannot be reduced to individual actor attributes. The whole exceeds the sum of parts.
\item \textbf{Self-organization produces structure.} Without central coordination, role differentiation and niche formation occur through local interactions.
\item \textbf{Co-evolution couples actors and context.} Role possibilities depend on what others do. Success by one actor reshapes opportunities for others.
\item \textbf{Non-linearity means surprises.} Small triggers can produce large effects (and vice versa). Tipping points and phase transitions are possible.
\end{itemize}

\subsection{Multimodal Identity Construction}

Organizational roles are not unitary---they are constructed differently across channels and audiences \cite{albert1985organizational,gioia2013organizational}. This dissertation demonstrates that identity is perspectival (what an organization ``is'' depends on who is looking and through what lens), gaps are strategic (divergence between claimed and enacted roles is positioning, not failure), and legitimacy is multimodal (organizations must construct coherent identities across channels, but strategic ambiguity can serve multiple audiences).

%===============================================================================
\section{Methodological Contributions}[Methods]
\labsec{methodological-contributions}
%===============================================================================

The dissertation advances methodology for studying innovation ecosystems in several ways.

\subsection{The RoleBox Framework}

\RoleBox{} provides an integrated computational infrastructure for multimodal role constellation analysis.\sidenote{\textbf{Open source}---All \RoleBox{} code is available at \texttt{github.com/babajideowoyele/rolebox} for replication and extension.} Its contributions include modular architecture (six data modality pipelines usable independently or in combination), reproducible workflows (documented pipelines enabling replication and adaptation), and open infrastructure (code, data schemas, and analytical procedures publicly available for critique, extension, and reuse).

\subsection{Multimodal Data Integration}

Combining five data modalities required developing integration protocols that align entity identities across sources, handle missing data, weight signals by reliability and relevance, and interpret divergence as signal rather than noise. These protocols transfer to other multimodal research contexts.

\subsection{Computational Role Extraction}

The dissertation demonstrates that NLP techniques can extract role-relevant information from unstructured text at scale:\sidenote{\textbf{Computational extraction}---NLP enables analysis at scales impossible manually. But human judgment remains essential for validation and interpretation.} role vocabulary identification through embedding-based semantic clustering, sociotype detection through community detection on role co-occurrence networks, and semantic role labeling for extracting actor-action-target triplets. These techniques complement rather than replace qualitative analysis.


%===============================================================================
\section{Bridging Data-Driven and Theory-Driven Research}[Bridging Paradigms]
\labsec{bridging-paradigms}
%===============================================================================

\textcite{maass2018data} identify fundamental challenges at the intersection of data-driven and theory-driven research. This dissertation's contribution lies precisely at this intersection: we bring \emph{multimodal organic data} (data-driven) to bear on \emph{role constellation theory} (theory-driven).

\subsection{Four Challenges}

Adapting \textcite{maass2018data} to sustainability transitions research, we identify four challenges specific to multimodal role constellation cartography:

\begin{description}[leftmargin=0.5cm, itemsep=0.3em]
\item[C1: Reconciling analytical approaches.] How can computational pattern discovery be reconciled with interpretive role theory?
\item[C2: Selecting data and techniques.] Which data sources and computational techniques are appropriate for which theoretical constructs?
\item[C3: Addressing previously intractable problems.] Can multimodal organic data answer questions about role dynamics that were impossible with designed data alone?
\item[C4: Enabling cumulative research.] Can role constellation data be shared, compared, and accumulated across research teams?
\end{description}

\subsection{Paths Between Perspectives}

Each perspective brings distinct strengths. Data-driven approaches offer scale, temporal coverage, multimodal breadth, and pattern discovery, but suffer from theoretical agnosticism and interpretation debt. Theory-driven approaches offer conceptual precision, interpretive depth, and causal reasoning, but face scale constraints and modality limits.\sidenote{\textbf{Complementary strengths}---Neither perspective alone is sufficient. Data-driven approaches lack theoretical grounding; theory-driven approaches lack scalability.}

The dissertation establishes explicit paths between data sources and theoretical constructs: websites map to sociotypes through NLP extraction, news maps to legitimacy through framing detection, social media maps to dynamics through network evolution, investment maps to functions through funding flows, and visual materials map to construction through image analysis.

\subsection{Resolution Through Cartographic Knowledge}

\cref{tab:challenge-resolution} summarizes how multimodal cartography resolves each challenge:

\begin{center}
\small
\tableminimal
\begin{tabular}{@{}p{0.5cm}p{3.5cm}p{4cm}p{4cm}@{}}
\toprule
& \minimalhead{Challenge} & \minimalhead{How Addressed} & \minimalhead{Key Artifacts} \\
\midrule
\textbf{C1} & Reconciling approaches & Explicit mapping between data sources and theoretical constructs & \cref{ch:method} \\
\addlinespace
\textbf{C2} & Selecting data/techniques & Theory-first selection: each modality operationalizes specific constructs & Design Principle DP1 \\
\addlinespace
\textbf{C3} & Addressing intractable problems & Scale (1000+ actors), temporal depth, multimodal triangulation & RoleBox pipelines \\
\addlinespace
\textbf{C4} & Enabling cumulative research & Open-source infrastructure, documented pipelines, shareable schemas & Design Principle DP7 \\
\bottomrule
\end{tabular}
\captionof{table}{Resolution of four challenges through multimodal cartography}
\labtab{challenge-resolution}
\end{center}

The dissertation fulfills three researcher roles identified by \textcite{maass2018data}: an \textbf{intermediary} role (translating role theory into computable features and interpreting clusters through sociotype frameworks), a \textbf{responsive} role (building \RoleBox{} to address the cartographic gap), and a \textbf{proactive} role (preparing open infrastructure for future researchers). These roles are not mutually exclusive; effective bridging research requires all three.

%===============================================================================
\section{Design Knowledge Contributions}[Design Knowledge]
\labsec{design-knowledge}
%===============================================================================

In Design Science Research \cite{hevner2004design,gregor2013positioning}, artifacts embody \textit{design knowledge}---transferable insights about \emph{how} to build effective solutions.\sidenote{\textbf{Design knowledge}---``Know-how'' that transfers beyond specific artifacts to inform future design efforts.} From the four Meta-Design Requirements and their instantiation in \RoleBox, we derive seven transferable design principles for multimodal research platforms:

\begin{center}
\small
\tableminimal
\begin{tabular}{@{}p{3.5cm}p{8.5cm}@{}}
\toprule
\minimalhead{Design Principle} & \minimalhead{Rationale and Application} \\
\midrule
\textbf{DP1: Theory-first data selection} & Let theoretical questions drive data collection, not data availability. For each modality, document what theoretical constructs it operationalizes. \\
\addlinespace
\textbf{DP2: Modular pipeline architecture} & Build independent processing pipelines for each data source. Enables gradual expansion, parallel development, and selective deployment. \\
\addlinespace
\textbf{DP3: Explicit entity resolution} & Treat identity alignment as a first-class problem. Same organizations appear differently across sources; invest in robust entity resolution early. \\
\addlinespace
\textbf{DP4: Divergence as signal} & Design integration to detect and preserve discrepancies between sources. Divergence reveals tensions, not just measurement error. \\
\addlinespace
\textbf{DP5: Human-in-the-loop validation} & Computational extraction requires human judgment for ambiguous cases. Build calibration, feedback, and correction mechanisms into workflows. \\
\addlinespace
\textbf{DP6: Layered abstraction} & Provide multiple access points: raw data for advanced users, processed features for analysts, visualizations for practitioners. \\
\addlinespace
\textbf{DP7: Reproducible documentation} & Document not just code but decisions, assumptions, and limitations. Future users inherit not just tools but understanding. \\
\bottomrule
\end{tabular}
\captionof{table}{Design principles for multimodal research platforms derived from \RoleBox development}
\labtab{design-principles}
\end{center}

Beyond principles, we identify recurring \textit{design patterns}---reusable solutions to common challenges:\sidenote{\textbf{Design patterns}---Proven solutions to recurring problems that can be adapted to new contexts.} multi-view triangulation (combining claimed, attributed, and enacted signals with divergence thresholds), temporal snapshot pipelines (re-runnable at intervals with timestamped deltas), semantic embedding integration (bridging vocabularies through shared latent dimensions), network-text co-analysis (combining topology with semantics), graduated automation (from manual annotation to incremental automation with human checkpoints), and LLM-based knowledge graph extraction (producing structured entity-relation-entity triples for rapid exploration).


%===============================================================================
\section{Practical Implications}[Implications]
\labsec{practical-implications}
%===============================================================================

What do these findings mean for practitioners---policymakers, intermediary organizations, and transition managers?

\subsection{For Transition Governance}

\begin{description}[leftmargin=0.5cm, itemsep=0.3em]
\item[Map before intervening.] Governance interventions should be informed by understanding existing role constellations.\sidenote{\textbf{Diagnostic cartography}---Before prescribing solutions, understand the existing landscape.} What functions are covered? Where are gaps?
\item[Design for emergence.] Top-down role assignment rarely works. Instead, create conditions that enable desired roles to emerge.
\item[Monitor constellation health.] Track not just individual actors but constellation-level properties: functional coverage, structural resilience, adaptive capacity.
\item[Expect non-linearity.] Policy interventions may have delayed, amplified, or counterintuitive effects.
\end{description}

\subsection{For Intermediary Organizations}

\begin{description}[leftmargin=0.5cm, itemsep=0.3em]
\item[Know your position.] Use multimodal analysis to understand how you are perceived across channels.\sidenote{\textbf{Self-awareness}---Organizations often have blind spots about how they are perceived externally.} Where do claimed and attributed roles diverge?
\item[Manage gaps strategically.] Some divergence is healthy---aspirational positioning creates room for growth. But persistent gaps between claims and enactment erode legitimacy.
\item[Differentiate or connect.] The two viable strategies in crowded ecosystems: specialize or bridge.
\end{description}

\subsection{For Innovation Ecosystems}

Ecosystems need diverse roles---not just entrepreneurs and investors, but also connectors, legitimators, capacity builders, and critics.\sidenote{\textbf{Role diversity}---Healthy ecosystems require the full complement of functions.} Having multiple actors capable of filling critical roles buffers against disruption. And ecosystem boundaries are permeable: actors participate in multiple constellations; resources and ideas flow across boundaries.

%===============================================================================
\section{Limitations}
\labsec{limitations}
%===============================================================================

This research has several important limitations that bound the generalizability of findings.\sidenote{\textbf{Honest boundaries}---Acknowledging limitations is not weakness but intellectual honesty.}

\begin{itemize}[leftmargin=*, itemsep=0.3em]
\item \textbf{Single ecosystem.} The EIT is one institutional configuration among many. Findings may not transfer to other innovation ecosystem models.
\item \textbf{European context.} EIT operates within European innovation policy. Role dynamics in other regions likely differ.
\item \textbf{Time period.} Data collection occurred 2020--2025. Role constellations evolve; these are snapshots, not eternal truths.
\item \textbf{Digital traces only.} Multimodal cartography captures what leaves digital footprints \cite{ruths2014social,boyd2012critical}. Offline interactions, tacit knowledge, informal networks remain invisible.
\item \textbf{English language bias.} Text analysis focused on English content. Non-English actors are underrepresented.
\item \textbf{Computational limitations.} NLP techniques introduce noise. Results should be triangulated with qualitative validation.
\item \textbf{Micro-meso focus.} The dissertation examines actor-level and constellation-level dynamics but does not fully address macro-level regime change.
\item \textbf{Descriptive orientation.} Normative questions (what \textit{should} constellations look like?) remain underexplored.
\end{itemize}

%===============================================================================
\section{Chapter Summary}
\labsec{synthesis-summary}
%===============================================================================

\begin{summarybox}
This chapter synthesized findings from the empirical studies and discussed their implications:

\begin{itemize}[leftmargin=*, itemsep=0.3em]
\item \textbf{Cross-modal triangulation} reveals that both convergence and divergence across data modalities are informative---convergence validates role attributions while divergence signals aspiration gaps, perception gaps, and strategic positioning.

\item \textbf{Cross-domain patterns} include persistent knowledge triangle imbalance, universal connector inflation, recurring sociotype configurations, systematic invisibility of peripheral actors, and \textbf{the breakdown gap}---complete absence of breakdown/exnovation roles across all five modalities.

\item \textbf{Theoretical contributions} include the CAS reconceptualization of role constellations, multimodal identity construction, and the cartographic metaphor for transitions research.

\item \textbf{Bridging paradigms}---adapting \textcite{maass2018data}, we articulate four challenges at the intersection of data-driven and theory-driven role constellation research, and demonstrate how cartographic knowledge resolves them.

\item \textbf{Design knowledge contributions} articulate seven design principles and six design patterns that transfer beyond this specific study to inform future multimodal research platform development.

\item \textbf{Practical implications} address transition governance (map before intervening, design for emergence), intermediary strategy (know your position, manage gaps), and ecosystem design (functional coverage, redundancy, permeability).

\item \textbf{Limitations} bound generalizability to specific empirical, methodological, and theoretical contexts.
\end{itemize}
\end{summarybox}

\vspace{1em}
Having synthesized findings and discussed their theoretical and practical implications, the concluding chapter articulates the dissertation's contributions, future research directions, and final reflections on multimodal cartography of role constellations.

\end{document}
